% !TeX root = ../main.tex
% !Mode:: "TeX:UTF-8"

\chapter{绪\quad 论}\label{chap:intro}

\section{研究背景与意义}

随着物联网（Internet of Things，IoT）与新一代无线通信技术的快速发展，智能终端和传感设备的数量呈爆发式增长。
在智慧城市、工业监控、环境感知、远程医疗等典型应用场景中，大量传感节点以一定的频率感知环境状态并产生状态更新
数据，这些数据经由无线网络汇聚至基站或边缘服务器进行处理。在此类实时监测系统中，信息能否在"失效"之前及时
送达接收端，直接决定了决策的准确性与系统的响应能力。例如，在自动驾驶中，若车辆感知到的前方障碍物信息过于陈旧，
则可能引发安全事故；在工业控制中，滞后的监测数据可能导致控制策略失效。

传统上，人们常用传输时延、吞吐量等指标衡量通信系统的性能，但这些指标难以刻画"信息的时效性"。
为此，\cite{2015-Pappas-ICC} 提出了信息新鲜度（Age of Information，AoI）这一概念，用于描述接收端所持有
信息相对于源端最新状态的新旧程度。AoI 综合衡量了状态生成、排队等待与传输时延三个环节，能够更准确地反映实时
监测场景中"信息是否足够新"这一关键需求。近年来，AoI 已成为通信与网络领域的研究热点，并被广泛用于物联网、
车联网、无人机通信等系统中\cite{2022-Jiao-TVT}。

另一方面，6G 物联网系统通常面临海量连接与有限资源的矛盾：可用频谱、发射功率与接收端处理能力均受约束，而
传感节点数量庞大、状态更新频繁。如何在有限的时频资源下，对多个数据源的状态更新进行合理调度，在保证系统队列
稳定的前提下最小化整体的信息陈旧程度，具有重要的理论意义与应用价值。同时，随着传感数据维度的不断提高，直接
传输原始高维数据会带来巨大的通信开销，如何利用机器学习方法对数据进行压缩与特征提取，也是值得研究的问题。

基于上述背景，本文以面向 6G 物联网的多源状态更新系统为对象，围绕"信息新鲜度"这一核心指标，研究数据采集与
传输调度的联合优化方法，为大规模实时物联网应用中的信息及时性保障提供理论依据与实现参考。

\section{国内外研究现状分析}

围绕多源传感网络中的数据采集与信息新鲜度保障，已有研究主要从信息新鲜度的度量与调度、队列稳定性优化以及
机器学习与网络资源管理的结合三个方面展开。

在信息新鲜度及其调度研究方面，\cite{2015-Pappas-ICC} 研究了具有队列管理的多源信息更新系统，分析了队列丢包
与排队策略对 AoI 性能的影响，指出合理的队列管理能够显著改善系统的信息时效性。\cite{2022-Jiao-TVT} 进一步
将信息新鲜度引入卫星互联网的双跳传输场景，提出了一种面向 AoI 优化的网络编码混合自动重传请求（HARQ）传输
方案，验证了 AoI 感知调度在时变信道下的有效性。上述工作表明，AoI 作为优化目标能够引导系统设计出时效性更好的
传输策略，但大多停留在单跳或特定信道模型下，对多源动态排队与联合调度的建模仍有待深入。

在队列稳定性与优化调度方面，李雅普诺夫优化（Lyapunov Optimization）提供了一套在不依赖未来统计信息的条件下
实现渐近最优控制的框架。通过在每时隙最小化"漂移加惩罚"项，可以在保证队列稳定的同时逼近最优的目标性能。这类
方法已被广泛用于无线网络中的吞吐量优化、能量管理与排队控制等场景，其本质是在"即时性能"与"队列稳定"之间取得
平衡，理论研究已较为成熟。

在机器学习与网络资源管理结合方面，堆叠自编码器（Stacked Autoencoder，SAE）作为一种无监督深度学习模型，
能够通过逐层预训练对高维输入数据进行非线性降维与特征提取，在数据压缩、异常检测、无线信号特征提取等任务中
表现出良好的性能。将其用于传感数据的压缩，可在保持关键信息的前提下显著降低传输比特数，为多源大规模采集提供
可能。

综合来看，现有研究在信息新鲜度度量、队列稳定调度以及数据压缩等方面分别取得了进展，但仍存在以下不足：一是对
多源状态更新与排队传输的联合系统建模不够完备；二是很少将深度学习特征压缩与 AoI 感知的在线调度相结合；三是
多数调度策略依赖未来信道或到达过程的统计分布，实用性受限。针对上述不足，本文提出一种"SAE 特征压缩 + 李雅普
诺夫漂移最小化调度"的结合方法，旨在克服这些问题。

\section{本文主要研究内容}

本文围绕面向 6G 物联网的多源状态更新系统，研究数据采集与传输调度的联合优化方法，主要研究内容如下：

\begin{enumerate}
  \item 建立了多源状态更新与排队传输的系统模型。考虑多个源节点经各自队列汇聚至基站的网络结构，以系统长时
        平均加权 AoI 为优化目标，在队列稳定约束下将其形式化为离散时间随机优化问题，明确了待求解的调度决策变量
        与约束条件。
  \item 提出了基于堆叠自编码器的状态数据压缩方法。针对高维传感数据的传输开销问题，利用 SAE 进行无监督逐层
        预训练，学习数据的低维特征表示，并给出其解码重建过程与损失函数，从而在降低传输负载的同时尽量保留状态
        信息。
  \item 设计了基于李雅普诺夫漂移最小化的在线调度算法。在每个时隙，依据当前各队列长度与 AoI 状态，通过最小化
        "漂移加惩罚"项确定最优的源节点进行传输，无需未来信道或到达统计信息，并证明了所提策略的渐近最优性与队列
        稳定性。
  \item 通过仿真实验验证了所提方法的有效性。在多源随机到达与动态信道条件下，对比了所提方法与轮询调度、最大
        AoI 优先、最大队列优先等基准算法的长时平均加权 AoI 与队列稳定性，并分析了 SAE 压缩率对整体时效性能的
        影响。
\end{enumerate}

\section{论文结构安排}

本文共分为六章，组织结构安排如下：

第 \cref{chap:intro} 章为绪论，介绍了研究背景与意义，综述了国内外在信息新鲜度、队列稳定调度与机器学习网络管理
方面的研究现状，并总结了本文的主要研究内容与章节安排。

第 2 章介绍相关理论基础。首先给出信息新鲜度（AoI）的基本定义与度量方法，然后阐述李雅普诺夫优化的基本原理，
接着介绍堆叠自编码器的结构与训练方法，最后给出本文所依赖的多源接入网络背景。

第 3 章建立系统模型并给出问题形式化。详细描述多源状态更新与排队传输的网络模型，推导 AoI 的演化方程，并在
队列稳定约束下将数据采集与调度问题建模为优化问题。

第 4 章提出本文的方法设计。分别给出基于堆叠自编码器的数据压缩方法、基于李雅普诺夫漂移最小化的在线调度算法，
并对算法的复杂度与理论性质进行分析。

第 5 章给出仿真实现与结果分析。介绍实验参数设置与数据集，从压缩性能与调度性能两个维度分析所提方法的优势，
并讨论关键参数的影响。

第 6 章总结全文工作，并对未来的研究方向进行展望。
